% ======================================================================
% SECTION 5: AUTONOMOUS SAFETY AND PHARMACOVIGILANCE
% File: sections/safety_pharmacovigilance.tex
% ======================================================================

Safety and pharmacovigilance represent the most consequential sponsor responsibility: the
obligation to protect trial participants from unreasonable risk while generating the safety
evidence required for regulatory decision-making. Under 21 CFR \S312.32, sponsors must
report serious and unexpected adverse reactions to the FDA within defined timelines
\cite{cfr-part-312}. ICH E2B(R3) mandates electronic Individual Case Safety Report (ICSR)
transmission, with the April 2026 deadline creating an immediate compliance requirement for
all sponsors \cite{ich-e2br3}. The autonomous safety system addresses these obligations
through a dedicated safety agent, an integrated safety management team interface, robotic
procedure safety monitoring (unique to Physical AI trials), and decentralized safety data
collection.

\subsection{Safety Agent}

The safety agent (safety\_agent) replaces the pharmacovigilance team for adverse event
intake, medical coding, causality assessment, signal detection, and expedited reporting. It
receives safety data from multiple sources: investigator site reports entered through the
EDC system, patient-reported outcomes collected through ePRO platforms, wearable device
telemetry transmitted through digital health technology integrations, and robotic procedure
telemetry from the robot execution gateway (\S\ref{sec:robotics}).

Adverse event intake and classification follows the MedDRA coding hierarchy. The agent
automatically codes reported events to preferred terms and system organ classes, identifies
seriousness criteria (death, life-threatening, hospitalization, disability, congenital
anomaly, or other medically important condition), and determines expectedness by comparing
the event against the Reference Safety Information in the current Investigator's Brochure.

Causality assessment implements three complementary methods. The WHO-UMC system provides
a structured categorical assessment (certain, probable, possible, unlikely, conditional,
unassessable) based on temporal relationship, pharmacological plausibility, dechallenge and
rechallenge evidence, and alternative explanations. The Naranjo algorithm provides a
numerical probability score through a ten-question standardized questionnaire. A Bayesian
causality model integrates prior probability from the known safety profile, likelihood
based on temporal and mechanistic evidence, and posterior probability computed through
Bayes' theorem. The three methods are applied in parallel, with discordant results
flagged for medical monitor review.

Signal detection uses statistical disproportionality analysis applied to the cumulative
adverse event database. The agent implements four complementary methods: Proportional
Reporting Ratio (PRR), Reporting Odds Ratio (ROR), Multi-item Gamma Poisson Shrinker
(MGPS), and Bayesian Confidence Propagation Neural Network (BCPNN). Signals are defined
as drug-event combinations where at least two of the four methods exceed their respective
thresholds. Detected signals enter a structured evaluation workflow: initial triage,
clinical review of individual cases, aggregate analysis, and determination of whether the
signal represents a confirmed safety risk requiring regulatory action.

Expedited reporting package generation automates compliance with 21 CFR \S312.32
\cite{cfr-part-312} and ICH E2A timelines. Fatal or life-threatening unexpected suspected
adverse reactions are reported within 7 calendar days, with follow-up within 8 additional
days. All other serious unexpected suspected adverse reactions are reported within 15
calendar days. The agent generates IND Safety Reports in the required format, routes them
for sponsor-of-record review (mandatory human gate for all expedited reports), and
transmits them electronically upon approval.

E2B(R3) ICSR transmission generates XML-formatted individual case safety reports
compliant with the ICH E2B(R3) standard \cite{ich-e2br3}. Each ICSR contains the
complete case narrative, MedDRA-coded events, reporter information, suspect and
concomitant medications, and relevant medical history. The agent validates each ICSR
against the E2B(R3) schema before transmission and maintains a complete audit trail
of all transmitted reports.

\subsection{Safety Management Team and DMC Interface}

The autonomous safety management team interface replaces the traditional safety review
board structure with an automated system that produces the same outputs: periodic safety
assessments, cumulative risk-benefit analysis, and safety-related recommendations for
study conduct modification. The system generates Development Safety Update Reports
(DSURs) per ICH E2F requirements \cite{ich-e2f}, producing annual comprehensive safety
analyses that include cumulative subject exposure, individual and aggregate adverse event
analysis, safety signal evaluation, and updated risk-benefit assessment.

The independent Data Safety Monitoring Committee interface provides the autonomous system
with a standardized mechanism for interacting with the external, independent DMC
required for controlled trials with serious outcomes. The agent prepares interim analysis
presentations in the format specified by the DMC charter, including efficacy and safety
data in unblinded or partially blinded formats as required. The DMC retains full authority
to recommend study modification, temporary suspension, or early termination. All DMC
recommendations are implemented through the study orchestrator with mandatory
sponsor-of-record notification.

The dose escalation committee interface supports oncology dose-finding studies using
established methods including the 3+3 design, modified toxicity probability interval
(mTPI) design, and Bayesian Optimal Interval (BOIN) design. The agent presents
dose-level safety data to the committee in standardized format, computes dose-limiting
toxicity rates, and implements committee dose decisions through the clinical
accountability agent. For Physical AI trials, the committee interface includes robotic
procedure safety data alongside pharmacological safety data, enabling integrated
dose and procedure risk assessment.

\subsection{Robotic Procedure Safety Monitoring}

Robotic procedure safety monitoring represents a capability unique to Physical AI trial
sponsors. Traditional safety monitoring addresses adverse drug reactions; the autonomous
sponsor must additionally monitor safety events arising from robotic clinical procedures,
creating a dual safety monitoring requirement that has no precedent in conventional
clinical trials.

The robotic safety event classification system categorizes events arising from Physical
AI procedures into four severity levels aligned with the adapted 21 CFR Part 312
\cite{cfr312-adapt} adverse event framework. Table~\ref{tab:robot-safety} presents
this classification.

\begin{longtable}{L{1.5cm} L{1.8cm} L{2.5cm} L{3.0cm} L{2.8cm}}
\caption{Robotic Safety Event Classification}
\label{tab:robot-safety} \\
\toprule
\textbf{Category} & \textbf{Severity} & \textbf{Detection Method} & \textbf{Response Protocol} & \textbf{Human Gate} \\
\midrule
\endfirsthead
\toprule
\textbf{Category} & \textbf{Severity} & \textbf{Detection Method} & \textbf{Response Protocol} & \textbf{Human Gate} \\
\midrule
\endhead
\midrule
\multicolumn{5}{r}{\textit{Continued on next page}} \\
\endfoot
\bottomrule
\endlastfoot
Category A: No deviation & Normal operation & Continuous telemetry within all parameters & Log and archive; no intervention & None required \\
Category B: Minor deviation & Low, self-correcting & Telemetry parameter outside normal range, auto-corrected within 5 seconds & Log deviation, continue procedure, increase monitoring frequency & None; automated logging \\
Category C: Significant deviation & Moderate, requires intervention & Telemetry parameter exceeds action threshold, not self-correcting & Pause procedure, notify supervising clinician, await instruction & Clinician approval required to resume \\
Category D: Critical event & High, potential patient harm & Emergency stop triggered, force limit exceeded, position deviation beyond safe envelope & Immediate procedure halt, secure patient, activate emergency protocol & Mandatory; no autonomous resumption \\
\end{longtable}

Robot telemetry integration ingests real-time data streams from surgical robots (da Vinci
Xi force/torque sensors, instrument position tracking, insufflation pressure), collaborative
robots (Franka Emika joint torque monitoring, end-effector force sensing, workspace boundary
compliance), and all other robotic systems operating within the trial site. The safety agent
correlates robotic telemetry with patient vital signs and clinical status to detect events
where robot operation may affect patient safety. This correlation analysis enables detection
of events that neither the robotic system nor the clinical monitoring system would identify
independently.

Table~\ref{tab:safety-timeline} presents the complete safety reporting timeline
requirements.

\begin{longtable}{L{2.2cm} L{2.2cm} L{1.8cm} L{2.8cm} L{2.5cm}}
\caption{Safety Reporting Timeline Requirements}
\label{tab:safety-timeline} \\
\toprule
\textbf{Report Type} & \textbf{Regulatory Basis} & \textbf{Timeline} & \textbf{Agent Responsibility} & \textbf{Output Format} \\
\midrule
\endfirsthead
\toprule
\textbf{Report Type} & \textbf{Regulatory Basis} & \textbf{Timeline} & \textbf{Agent Responsibility} & \textbf{Output Format} \\
\midrule
\endhead
\midrule
\multicolumn{5}{r}{\textit{Continued on next page}} \\
\endfoot
\bottomrule
\endlastfoot
IND Safety Report (fatal or life-threatening) & 21 CFR \S312.32(c)(1) & 7 calendar days (initial) & safety\_agent with mandatory sponsor-of-record gate & FDA Form 3500A with narrative \\
IND Safety Report (serious, unexpected) & 21 CFR \S312.32(c)(1) & 15 calendar days & safety\_agent with mandatory sponsor-of-record gate & FDA Form 3500A with narrative \\
Follow-up IND Safety Report & 21 CFR \S312.32(d) & 15 calendar days from new information & safety\_agent & Amended FDA Form 3500A \\
E2B(R3) ICSR & ICH E2B(R3) & Per jurisdiction (typically 15 days) & safety\_agent & E2B(R3) XML \\
DSUR & ICH E2F & Annual (within 60 days of DIBD) & safety\_agent & Formatted DSUR document \\
IND Annual Report & 21 CFR \S312.33 & Within 60 days of anniversary & regulatory\_agent & Formatted annual report \\
Robotic Safety Event Report & Adapted 21 CFR 312 & Category C: 24 hours; Category D: immediate & robot\_execution\_gateway & Structured event report with telemetry \\
\end{longtable}

\subsection{Decentralized Safety Monitoring}

Decentralized safety monitoring extends the safety data collection beyond the investigator
site through digital health technologies, wearable devices, and patient-reported outcome
instruments. This capability is essential for Physical AI trials that incorporate
decentralized elements such as remote follow-up visits, home-based rehabilitation with
exoskeleton robots, and remote monitoring between in-person treatment visits.

Digital health technology validation follows FDA guidance on digital health technologies for
remote data acquisition. The safety agent validates that each DHT device meets the data
quality requirements specified in the protocol, including measurement accuracy, data
transmission reliability, and patient compliance monitoring. Wearable device data (heart
rate, activity level, sleep patterns, temperature) is integrated into the safety monitoring
stream with appropriate clinical context to enable detection of safety signals that might
not be captured through periodic site visits.

The integration of decentralized safety data with site-reported data creates a comprehensive
safety monitoring environment that exceeds the capabilities of traditional pharmacovigilance.
Traditional safety monitoring relies primarily on investigator-reported adverse events, which
are subject to reporting delays (the time between event occurrence and investigator awareness),
assessment variability (different investigators may classify the same event differently), and
recall bias (patients may not report symptoms accurately at periodic visits). The autonomous
safety system supplements investigator reporting with continuous data streams from wearable
devices, ePRO instruments, and robotic telemetry, providing a multi-modal safety surveillance
capability that detects events earlier and classifies them more consistently.

The safety agent maintains a cumulative safety database that integrates all safety data
sources into a unified repository. This repository enables longitudinal safety analysis across
the trial lifecycle: from the first patient enrolled through database lock and beyond into
post-marketing surveillance. The database supports the aggregate safety analyses required for
DSURs, the signal detection analyses that run continuously during the trial, and the
individual case narratives required for regulatory submissions. Data quality controls ensure
that all entries are validated, coded to current MedDRA and WHODrug dictionaries, and linked
to their source documentation.

For Physical AI trials specifically, the safety agent must address the intersection of drug
safety and device safety. When a patient experiences an adverse event during a robotic
procedure, the safety agent must determine whether the event is related to the investigational
drug, the robotic procedure, their interaction, the underlying disease, or other factors. This
multi-factorial causality assessment extends beyond traditional pharmacovigilance methods and
requires integration of pharmacological knowledge, device performance data, and procedural
context. The safety agent implements a structured decision tree that evaluates each potential
causal factor and documents the reasoning behind the causality determination.

Patient-reported outcomes for safety use validated electronic instruments deployed through
ePRO platforms. The agent configures PRO-CTCAE (Patient-Reported Outcomes version of the
Common Terminology Criteria for Adverse Events) and other validated instruments appropriate
for the trial indication and treatment modality. Responses are analyzed in real time with
automated flagging of symptom severity increases that warrant clinical follow-up. This
approach provides continuous patient safety monitoring between scheduled assessment visits,
reducing the risk of undetected adverse events in the interval between site contacts.
